\clearpage
\section{Systemkontext}

Der fachliche Kontext umfasst vier Beteiligte: den Disponenten, DISPO, das System zur
digitalen Lieferavisierung und den Kunden. Der Disponent plant in DISPO und löst dort die
Kommunikation aus. Der Kunde antwortet über die Bestätigungsseite. Das Ergebnis wird dem
Disponenten im Dashboard und perspektivisch auch in DISPO angezeigt.

\begin{figure}[H]
\centering
\resizebox{0.98\textwidth}{!}{%
\begin{tikzpicture}[font=\scriptsize,>=Latex]
  \node[draw=MinovaDark,rounded corners=2pt,minimum width=2.4cm,minimum height=1.3cm,align=center,fill=SoftGray]
    (dispatcher) at (0,0) {\textbf{Disponent}\\plant und entscheidet};
  \node[draw=MinovaDark,rounded corners=2pt,minimum width=2.5cm,minimum height=1.4cm,align=center,fill=white]
    (dispo) at (3.5,0) {\textbf{DISPO}\\führendes System};
  \node[draw=MinovaBlue,very thick,rounded corners=3pt,minimum width=3.8cm,minimum height=2.4cm,align=center,fill=SoftBlue]
    (system) at (8.5,0) {\textbf{Digitale Lieferavisierung}\\[2pt]Avisierungsservice\\Bestätigungsseite\\Avisierungsdashboard};
  \node[draw=MinovaDark,rounded corners=2pt,minimum width=2.4cm,minimum height=1.3cm,align=center,fill=SoftGray]
    (customer) at (13.6,0) {\textbf{Kunde}\\bestätigt oder lehnt ab};

  \draw[->,thick] (dispatcher) -- (dispo);
  \draw[<->,thick] (dispo) -- node[above=5pt,font=\tiny,fill=white,inner sep=1pt]{Daten / Ergebnis} (system);
  \draw[<->,thick] (system) -- node[above=5pt,font=\tiny,fill=white,inner sep=1pt]{Nachricht / Antwort} (customer);
\end{tikzpicture}
}
\caption{Fachlicher Systemkontext der digitalen Lieferavisierung}
\label{fig:system-context}
\end{figure}


\begin{center}
\begin{tabularx}{0.94\textwidth}{>{\bfseries}L{3.15cm}X}
\hline
Disponent & Löst die Avisierung aus und bearbeitet Fälle mit Handlungsbedarf.\\
\hline
DISPO & Liefert die geplanten Daten und übernimmt die Ergebnisanzeige.\\
\hline
Lieferavisierung & Versendet die Nachricht, erfasst die Antwort und bereitet sie auf.\\
\hline
Kunde & Prüft das vorgeschlagene Zeitfenster und antwortet über die Web-Seite.\\
\hline
\end{tabularx}
\end{center}

Details des eigentlichen Nachrichtenversands sind in dieser Übersicht bewusst zu einem
allgemeinen Kommunikationskanal zusammengefasst. Für das Verständnis des Gesamtablaufs ist
nur entscheidend, dass die Nachricht den persönlichen Link zur Bestätigungsseite enthält.
